\documentclass[10pt,a4paper]{amsart}
\usepackage[margin=1in]{geometry}
\usepackage{barnette-figs}
\usepackage{graphicx,amssymb}
\setlength{\parindent}{0pt}
\setlength{\parskip}{7pt}
\pagestyle{plain}
\begin{document}
\section*{17. What a small cut must remember}
\begin{center}\begin{minipage}{.86\textwidth}\centering
\begin{tikzpicture}[bf,x=1cm,y=1cm]
\draw[bf edge] (0.590,0.000) -- (-0.295,0.511);
\draw[bf edge] (-0.295,0.511) -- (-0.295,-0.511);
\draw[bf edge] (-0.295,-0.511) -- (0.590,0.000);
\node[draw,fill=black,circle,minimum size=4pt,inner sep=0pt] at (0.590,0.000) {};
\node[draw,fill=black,circle,minimum size=4pt,inner sep=0pt] at (-0.295,0.511) {};
\node[draw,fill=white,circle,minimum size=4pt,inner sep=0pt] at (-0.295,-0.511) {};
\draw[bf edge] (2.340,0.000) -- (1.455,0.511);
\draw[bf edge] (1.455,0.511) -- (1.455,-0.511);
\draw[bf edge] (1.455,-0.511) -- (2.340,0.000);
\node[draw,fill=black,circle,minimum size=4pt,inner sep=0pt] at (2.340,0.000) {};
\node[draw,fill=white,circle,minimum size=4pt,inner sep=0pt] at (1.455,0.511) {};
\node[draw,fill=black,circle,minimum size=4pt,inner sep=0pt] at (1.455,-0.511) {};
\draw[bf edge] (4.090,0.000) -- (3.205,0.511);
\draw[bf edge] (3.205,0.511) -- (3.205,-0.511);
\draw[bf edge] (3.205,-0.511) -- (4.090,0.000);
\node[draw,fill=black,circle,minimum size=4pt,inner sep=0pt] at (4.090,0.000) {};
\node[draw,fill=white,circle,minimum size=4pt,inner sep=0pt] at (3.205,0.511) {};
\node[draw,fill=white,circle,minimum size=4pt,inner sep=0pt] at (3.205,-0.511) {};
\draw[bf edge] (0.590,-2.150) -- (0.000,-1.560);
\draw[bf edge] (0.000,-1.560) -- (-0.590,-2.150);
\draw[bf edge] (-0.590,-2.150) -- (-0.000,-2.740);
\draw[bf edge] (-0.000,-2.740) -- (0.590,-2.150);
\node[draw,fill=black,circle,minimum size=4pt,inner sep=0pt] at (0.590,-2.150) {};
\node[draw,fill=black,circle,minimum size=4pt,inner sep=0pt] at (0.000,-1.560) {};
\node[draw,fill=black,circle,minimum size=4pt,inner sep=0pt] at (-0.590,-2.150) {};
\node[draw,fill=white,circle,minimum size=4pt,inner sep=0pt] at (-0.000,-2.740) {};
\draw[bf edge] (2.340,-2.150) -- (1.750,-1.560);
\draw[bf edge] (1.750,-1.560) -- (1.160,-2.150);
\draw[bf edge] (1.160,-2.150) -- (1.750,-2.740);
\draw[bf edge] (1.750,-2.740) -- (2.340,-2.150);
\node[draw,fill=black,circle,minimum size=4pt,inner sep=0pt] at (2.340,-2.150) {};
\node[draw,fill=black,circle,minimum size=4pt,inner sep=0pt] at (1.750,-1.560) {};
\node[draw,fill=white,circle,minimum size=4pt,inner sep=0pt] at (1.160,-2.150) {};
\node[draw,fill=black,circle,minimum size=4pt,inner sep=0pt] at (1.750,-2.740) {};
\draw[bf edge] (4.090,-2.150) -- (3.500,-1.560);
\draw[bf edge] (3.500,-1.560) -- (2.910,-2.150);
\draw[bf edge] (2.910,-2.150) -- (3.500,-2.740);
\draw[bf edge] (3.500,-2.740) -- (4.090,-2.150);
\node[draw,fill=black,circle,minimum size=4pt,inner sep=0pt] at (4.090,-2.150) {};
\node[draw,fill=white,circle,minimum size=4pt,inner sep=0pt] at (3.500,-1.560) {};
\node[draw,fill=black,circle,minimum size=4pt,inner sep=0pt] at (2.910,-2.150) {};
\node[draw,fill=black,circle,minimum size=4pt,inner sep=0pt] at (3.500,-2.740) {};
\draw[bf edge] (5.840,-2.150) -- (5.250,-1.560);
\draw[bf edge] (5.250,-1.560) -- (4.660,-2.150);
\draw[bf edge] (4.660,-2.150) -- (5.250,-2.740);
\draw[bf edge] (5.250,-2.740) -- (5.840,-2.150);
\node[draw,fill=black,circle,minimum size=4pt,inner sep=0pt] at (5.840,-2.150) {};
\node[draw,fill=white,circle,minimum size=4pt,inner sep=0pt] at (5.250,-1.560) {};
\node[draw,fill=black,circle,minimum size=4pt,inner sep=0pt] at (4.660,-2.150) {};
\node[draw,fill=white,circle,minimum size=4pt,inner sep=0pt] at (5.250,-2.740) {};
\draw[bf edge] (7.590,-2.150) -- (7.000,-1.560);
\draw[bf edge] (7.000,-1.560) -- (6.410,-2.150);
\draw[bf edge] (6.410,-2.150) -- (7.000,-2.740);
\draw[bf edge] (7.000,-2.740) -- (7.590,-2.150);
\node[draw,fill=black,circle,minimum size=4pt,inner sep=0pt] at (7.590,-2.150) {};
\node[draw,fill=white,circle,minimum size=4pt,inner sep=0pt] at (7.000,-1.560) {};
\node[draw,fill=white,circle,minimum size=4pt,inner sep=0pt] at (6.410,-2.150) {};
\node[draw,fill=white,circle,minimum size=4pt,inner sep=0pt] at (7.000,-2.740) {};
\end{tikzpicture}
\end{minipage}\end{center}
The panels restrict every Hamiltonian partition of the cube's octahedral dual to one triangular and one quadrilateral boundary. Labels are identified under swapping the two tree classes. The three triangle patterns are exactly the choices of the unique monochromatic boundary edge. Across a separating triangle, the pieces can therefore agree by prescribing that edge. Quadrilateral boundaries carry more information; the five displayed label words are computed restrictions in this example. They are not the six piece states of the report, which also record boundary connectivity. A general completeness statement uses the stronger path-avoidance consequence of the supplied proof, and the categorical report still leaves one subcase to finite checking. At length five, the report finds that the realized patterns depend on the piece. Thus gluing cannot be replaced by a single unrestricted boundary rule for all cut lengths. The atlas distinguishes these observed restrictions from a new inductive existence proof.

\textbf{Status.} COMPUTED cube boundary restrictions here. LIT triangular gluing. Categorical 4-cut completeness proof is unrefereed and has a finite-check subcase; 5-cut failure is unrefereed computation.

\textbf{Connects to.} 18 cut side; 19 recovery; 24 earlier decomposition methods.

{\footnotesize\raggedright Sources: papers/\allowbreak{}OpenAI-Barnette.pdf, Section 7;\allowbreak{} research/\allowbreak{}2026-11-14\_fable\_route-categorical/\allowbreak{}REPORT.md;\allowbreak{} computed.json, cube\_boundary.\par}

\end{document}